\documentclass{article}
\title{Practice Problem 2.35}
\author{CSAPP}
\usepackage{booktabs, amssymb, amsmath, listings, parskip, mathtools, solarized-light}
\begin{document}
\maketitle

\begin{flushleft}
    You are given the assignment to develop code for a function $tmult\_ok$ that will
    determine whether two arguments can be multiplied without causing overflow.
\small
\begin{lstlisting}[language=c]
/* Determine whether arguments can be multiplied without overflow */
#include <stdio.h>
int tmult_ok(int x, int y){
    int p = x * y;
    /* Either x is zero, or dividing p by x gives y */
    return !x || p/x == y;
}
\end{lstlisting}

Devise a mathematical justification of your approach, along the following lines.
First, argue that the case $x = 0$ is handled correctly. Otherwise, consider $w$-bit
numbers $x\ (x \neq 0),\ y,\ p,$ and $q$, where $p$ is the result of performing two's complement
multiplication on $x$ and $y$, and $q$ is the result of dividing $p$ by $x$.

\bigskip
$\mathbf{1.}$ Show that $x \cdot y$, the integer product of $x$ and $y$, can be written in the form
$x \cdot y = p + t2^w$, where $t \neq 0$ if and only if the computation of $p$ overflows.

$\mathbf{2.}$ Show that $p$ can be written in the form $p = x \cdot q + r$, where $|r| < |x|$.

$\mathbf{3.}$ Show that $q = y$ if and only if $r = t = 0$.

\bigskip
\textbf{\normalsize SOLUTION}

\bigskip
\textbf{Case $x = 0$:}\\
If $x = 0$, then \texttt{!x} evaluates to true and the function returns $1$, indicating no overflow.
This is correct because $0 \cdot y = 0$ for any $y$, which never overflows. \hfill$\square$

\bigskip
\textbf{Part 1:}\\
Two's complement multiplication computes the lower $w$ bits of the true integer product $x \cdot y$.
Formally, there exists an integer $t$ such that:
\begin{equation}
    x \cdot y = p + t2^w
\end{equation}
where $p$ is the $w$-bit two's complement result. If no overflow occurs, $p$ exactly represents
$x \cdot y$, so $t = 0$. If overflow occurs, the upper bits are truncated and $t \neq 0$ accounts
for the discarded value. Therefore:
\begin{equation*}
    t \neq 0 \iff \text{overflow occurs} \hfill\square
\end{equation*}

\bigskip
\textbf{Part 2:}\\
By the division algorithm, for any integers $p$ and $x$ with $x \neq 0$, there exist unique
integers $q$ and $r$ satisfying:
\begin{equation}
    p = x \cdot q + r, \quad |r| < |x|
\end{equation}
In C, integer division \texttt{p/x} computes exactly this truncating quotient $q$,
with remainder $r = p - x \cdot q$. \hfill$\square$

\bigskip
\textbf{Part 3:}\\
From Part 1 we have $p = x \cdot y - t2^w$. Substituting into Part 2:
\begin{align}
    x \cdot q + r &= x \cdot y - t2^w \\
    x(y - q)      &= r + t2^w
\end{align}

\textit{($\Leftarrow$)} Suppose $r = t = 0$. Then $x(y - q) = 0$. Since $x \neq 0$, we conclude $q = y$.

\medskip
\textit{($\Rightarrow$)} Suppose $q = y$. Then:
\begin{equation}
    x(y - y) = 0 = r + t2^w \implies r = -t2^w
\end{equation}
From Part 2, $|r| < |x| \leq 2^{w-1}$. But if $t \neq 0$, then $|t2^w| \geq 2^w > 2^{w-1} > |r|$,
a contradiction. Therefore $t = 0$, which forces $r = 0$. \hfill$\square$

\bigskip
\textbf{Conclusion:}\\
The check \texttt{p/x == y} tests whether $q = y$, which by Part 3 holds if and only if
$t = 0$, i.e., no overflow occurred. Combined with the guard for $x = 0$, the function
\texttt{tmult\_ok} correctly detects two's complement multiplication overflow in all cases.

\end{flushleft}
\end{document}